\section{Mathematical Definition and Properties}
\label{sec:definition}

\subsection{The Seats-Votes Function}

The Partisan Bias is defined as a scalar summary of the seats-votes function $S(v)$.
For a $k$-district congressional plan with district-level Democratic vote shares
$v_1, v_2, \ldots, v_k$ (observed for a specific election cycle), the uniform
swing model estimates the seats-votes function as follows.

Let $\bar{v}_D = k^{-1}\sum_i v_i$ be the observed mean Democratic vote share.
The uniform swing needed to bring the statewide mean to a target value $v$ is
$\delta(v) = v - \bar{v}_D$. Applying this swing to all districts:
\[
\hat{v}_i(v) = v_i + \delta(v) = v_i + v - \bar{v}_D.
\]
The estimated Democratic seat share at statewide vote share $v$ is:
\[
S(v) = \frac{1}{k}\sum_{i=1}^k \mathbf{1}[\hat{v}_i(v) > 0.50].
\]
The function $S(v)$ is a step function that jumps at values of $v$ where a
swing brings a district exactly to 50\% ($v = 0.50 - v_i + \bar{v}_D$).

\subsection{Partisan Bias Definition}

Partisan Bias is:
\[
\text{Bias} = S(0.50) - 0.50.
\]
$\text{Bias} > 0$ indicates Democratic advantage at the 50\% pivot (Democrats
would win more than half the seats with exactly half the votes). $\text{Bias} < 0$
indicates Republican advantage. Because $S(v)$ is a step function, Bias is computed
at $v = 0.50$, using the convention that ties are broken by the observed order of
district vote shares.

For NC ($k=14$), $S(v)$ takes values $\{0, 1/14, 2/14, \ldots, 14/14\}$.
The minimum nonzero $|\text{Bias}|$ is $1/14 \approx 0.071$, corresponding
to a one-seat advantage. Enacted NC Bias $\approx -0.07$ thus indicates that
Republicans would win exactly one additional seat (8 vs.\ 7) if the statewide
vote were tied---a systematic structural advantage equivalent to one seat in a
14-seat delegation.

\subsection{The Symmetry Standard}

Zero Partisan Bias is a necessary condition for the full symmetry standard but
not sufficient. The full symmetry standard requires:
\[
S(v) = 1 - S(1-v) \quad \text{for all } v \in [0,1].
\]
This means the seats-votes function is anti-symmetric around the point $(0.50, 0.50)$:
the curve is the same whether read from left to right or reflected. Equivalently,
the party winning 60\% of votes should win the same seat share that the opposing
party would win at 60\% of votes.

Partisan Bias at the 50\% pivot tests only whether $S(0.50) = 0.50$, not whether
the full symmetry condition holds across all vote levels. A map could satisfy the
Bias test at 50\% while having asymmetric responsiveness at 55\% and 45\%. The
seats-votes curve (L.5) provides the full symmetry picture; Partisan Bias provides
a scalar summary at the most policy-relevant vote level.

\subsection{Uniform Swing Limitation}

The uniform swing model applies the same $\delta$ to every district, implicitly
assuming that all districts respond identically to statewide vote changes. This
assumption fails systematically in two scenarios:

\textbf{Safe-seat heterogeneity.} Districts with vote shares far from 50\% (safe
seats) respond less to statewide swings in practice, because incumbency advantage,
uncontested races, and campaign resource allocation protect them from the full
swing magnitude. Applying uniform swing to safe seats overestimates how many
seats change hands, which overestimates the slope of $S(v)$ near 50\% and
biases the Partisan Bias estimate.

\textbf{Geographic clustering.} Urban districts respond differently from suburban
and rural districts because the vote-share determinants (income, education, race)
change nonlinearly with geography. Uniform swing treats all districts as exchangeable,
ignoring this geographic heterogeneity.

For Texas ($k=38$), which has approximately 12 safe Republican districts ($v_i < 0.35$)
and 8 safe Democratic districts ($v_i > 0.65$), uniform swing overestimates
$|\text{Bias}|$ by approximately 15\% compared to a heteroskedastic correction
that applies reduced swing magnitudes to safe districts. The qualitative conclusion
(enacted TX $|\text{Bias}| \approx 0.09$ vs.\ bisect $|\text{Bias}| < 0.02$)
is robust to this correction.

\subsection{Swing Ratio (Electoral Responsiveness)}
\label{ssec:swing-ratio}

\begin{definition}[Swing Ratio]
\label{def:swing-ratio}
The \emph{swing ratio} (also called \emph{electoral responsiveness}) is the slope
of the seats-votes curve $S(v)$ at the statewide vote share $v = 0.50$:
\[
  \mathrm{SR} \;=\; \left.\frac{dS}{dv}\right|_{v=0.50}.
\]
We follow the Gelman-King (1994) definition~\citep{gelman1994}: SR is estimated
by a linear regression of seat share on vote share across a set of historical
elections, evaluated at $v = 0.50$.
\end{definition}

\paragraph{Why Gelman-King over Tufte (1973).}
Tufte defined the swing ratio as the point-slope at a specific election year, not
the regression-based estimate. McDonald-Burnett (2022) introduced a bounded version
robust to outlier elections. We use Gelman-King because: (1) it is the standard
in redistricting litigation expert reports; (2) it produces confidence intervals
via standard regression inference; (3) it is the definition referenced in McGhee
(2014) and Stephanopoulos-McGhee (2015), which are the primary efficiency gap
papers in our program.

A SR of 2 means a 1-percentage-point swing in votes produces a 2-percentage-point
swing in seats — the historical U.S.\ norm under single-member plurality elections.
A SR below 2 indicates a ``safe-seat'' map where vote changes have muted seat effects
(incumbency protection or packing). A SR above 2 indicates a ``responsive'' map where
small vote changes produce large seat changes.

\paragraph{Bisect swing ratios.}
For \textsc{bisect} maps across NC, WI, and TX: SR $\approx 2.1$--$2.3$, close to
the national historical average, consistent with compact maps not having unusual
safe-seat concentrations. Enacted maps in NC and TX have SR $\approx 1.5$--$1.8$,
reflecting the safe-seat buffer that protects incumbents of the mapmaking party.
